% ==================================================
% Section 6
% ==================================================
\section{Sampling-Time Effect and a Bilinear (Tustin) Redesign}

The discrete-time controller of Section~5 used the simplest approximations: the
error derivatives were computed with backward differences and the adaptation law
was integrated with forward Euler. These are only accurate when the sampling
period \(T_s\) is small. In this section we first show how the closed-loop
performance gets worse as \(T_s\) grows, explain why, and then present an
improved discrete-time design based on the bilinear (Tustin) transform.

\subsection{Performance Degradation with Increasing Sampling Time}

The Section~5 controller was simulated for several sampling periods, keeping the
continuous-time plant, the reference model and all controller gains the same.
Figure~\ref{fig:s6_degradation} shows the tracking error \(e(t)\) and the
filtered error \(z(t)\) for \(T_s = 0.01,\,0.04,\,0.05~\mathrm{s}\). The peak
tracking error grows as \(T_s\) increases, and the filtered error \(z(t)\) shows
a larger spike at start-up.

\begin{figure}[H]
\centering
\includegraphics[width=0.7\textwidth]{figs/hw/section6_degradation.png}
\caption{Baseline (Section~5) controller: tracking error \(e(t)\) (top) and
filtered error \(z(t)\) (bottom) for increasing sampling periods.}
\label{fig:s6_degradation}
\end{figure}

\noindent The main cause is the second backward difference used to estimate
\(\ddot{e}\),
\begin{equation}
\ddot{e}[k] \approx \frac{e[k]-2e[k-1]+e[k-2]}{T_s^{2}},
\label{eq:s6_secdiff}
\end{equation}
whose gain grows like \(1/T_s^{2}\). Together with the half-sample delay of the
zero-order hold, this term amplifies the sampled error and feeds it back through
\(z[k]=\ddot{e}[k]+\lambda_2\dot{e}[k]+\lambda_1 e[k]\) and the control law. As
\(T_s\) grows the loop loses phase margin, and at \(T_s\approx0.06~\mathrm{s}\)
the closed loop becomes unstable and the simulation diverges
(Table~\ref{tab:s6_compare}). The forward-Euler integration of the adaptation law
adds a smaller error of the same kind.

\subsection{Improved Design: Bilinear (Tustin) Discretisation}

To recover good performance without reducing \(T_s\), the controller is
re-discretised with the bilinear (Tustin) transform,
\begin{equation}
s \;\longleftarrow\; \frac{2}{T_s}\,\frac{z-1}{z+1},
\label{eq:s6_tustin}
\end{equation}
which gives a more accurate approximation of the continuous-time dynamics than
the forward-difference rule \(s\leftarrow(z-1)/T_s\).

\paragraph{Filtered (dirty) derivatives.}
A pure differentiator cannot be implemented directly, so each error derivative is
taken from a first-order filtered (``dirty'') derivative
\(H(s)=\dfrac{s}{\tau s+1}\), which is proper and band-limited (\(\tau\) is a
small filter time constant). Applying the bilinear
transform~\eqref{eq:s6_tustin} to \(H(s)\) and writing \(a = 2/T_s\) gives
\begin{equation}
\dot{\hat{e}}[k]=
\frac{a\big(e[k]-e[k-1]\big)-(1-\tau a)\,\dot{\hat{e}}[k-1]}{\tau a+1},
\label{eq:s6_dirty}
\end{equation}
and \(\ddot{\hat{e}}[k]\) is obtained by applying the same filter again to
\(\dot{\hat{e}}[k]\). Unlike~\eqref{eq:s6_secdiff}, the gain of
\eqref{eq:s6_dirty} is bounded, so high-frequency content is not amplified by
\(1/T_s^{2}\).

\paragraph{Trapezoidal adaptation.}
The adaptation law \(\dot{\hat{\theta}}=-\Gamma\,\varphi z\) is integrated with
the trapezoidal (Tustin) rule instead of forward Euler:
\begin{equation}
\hat{\theta}[k]=\hat{\theta}[k-1]
-\frac{T_s}{2}\,\Gamma\big(\varphi[k]z[k]+\varphi[k-1]z[k-1]\big).
\label{eq:s6_trap}
\end{equation}
The filtered error \(z[k]\), the auxiliary reference \(y_r^{(3)}[k]\) and the
control law \(u[k]\) are computed as before, using \eqref{eq:s6_dirty} for the
derivatives.

\subsection{Simulation Results and Discussion}

Table~\ref{tab:s6_compare} compares the peak tracking error of the baseline and
the redesigned controller for different \(T_s\) (\(\tau=0.06~\mathrm{s}\)).

\begin{table}[H]
\centering
\begin{tabular}{c c c}
\hline
\(T_s\) [s] & Baseline \(\max|e|\) & Tustin \(\max|e|\)\\
\hline
0.01  & 0.038 & 0.064 \\
0.04  & 0.085 & 0.070 \\
0.05  & 0.105 & 0.082 \\
0.06  & \emph{diverges} & 0.104 \\
0.07  & \emph{diverges} & 0.090 \\
0.075 & \emph{diverges} & 0.105 \\
0.08  & \emph{diverges} & \emph{diverges}\\
\hline
\end{tabular}
\caption{Peak tracking error \(\max_{t>1\,\mathrm{s}}|e(t)|\) versus sampling
period.}
\label{tab:s6_compare}
\end{table}

\noindent Two things can be seen. First, at the largest sampling period where the
baseline is still stable (\(T_s=0.05~\mathrm{s}\)) the redesign lowers the peak
error from \(0.105\) to \(0.082~\mathrm{rad}\) and reduces the overshoot, as
shown in Figure~\ref{fig:s6_fix}. Second, the redesign stays stable up to
\(T_s\approx0.075~\mathrm{s}\), well past the \(T_s\approx0.06~\mathrm{s}\) where
the baseline diverges, so the microcontroller can run at a lower sampling rate.

\begin{figure}[H]
\centering
\includegraphics[width=0.7\textwidth]{figs/hw/section6_tustin_fix.png}
\caption{Closed loop at \(T_s=0.05~\mathrm{s}\): output tracking (top) and
tracking error (bottom), for the baseline and the Tustin redesign.}
\label{fig:s6_fix}
\end{figure}

\noindent Figure~\ref{fig:s6_params} shows that the parameter estimates of the
redesigned controller at \(T_s=0.06~\mathrm{s}\) -- where the baseline cannot run
-- stay bounded and follow the same trajectories as the well-sampled reference
(\(T_s=0.01~\mathrm{s}\)). As in Sections~4 and~5 they do not converge to the
true values, because the step reference is not persistently exciting.

\begin{figure}[H]
\centering
\includegraphics[width=0.7\textwidth]{figs/hw/section6_parameter_comparison.png}
\caption{Parameter estimates \(\hat{\theta}_i(t)\): well-sampled baseline
(\(T_s=0.01~\mathrm{s}\)) and the Tustin redesign at \(T_s=0.06~\mathrm{s}\).}
\label{fig:s6_params}
\end{figure}

\noindent The only downside of the redesign is a slightly larger error at very
small sampling periods (\(T_s=0.01~\mathrm{s}\) in Table~\ref{tab:s6_compare}),
where the dirty-derivative filter adds a small phase lag that the raw backward
difference does not. This is a small price, since the point of the redesign is to
allow coarse sampling, and in that range it is consistently better. In short,
discretising the controller with the bilinear transform -- filtered
derivatives~\eqref{eq:s6_dirty} and trapezoidal integration of the adaptation
law~\eqref{eq:s6_trap} -- reduces the performance loss caused by the
discrete-time implementation and allows a larger sampling period.

\bigskip
